% Source: Wald GR, physical PDF pp. 120--129 (printed pp. 107--116).
% Coverage: Section 5.4, The Evolution of Our Universe; equations (5.4.1)--(5.4.11).
% Figures/tables/footnotes: no figures or tables; footnotes 1--2.
% Status: source-reviewed and compile-clean.
% Uncertainties: none.

\subsection{The Evolution of Our Universe}\label{sec:5.4}

In this section we shall briefly outline the history of our universe from the big bang
to the present. The picture we shall present is the ``standard'' one, which assumes that
the universe is well described throughout its history by a Robertson--Walker solution
and which makes further assumptions concerning the matter content of the universe.
While this picture has received strong supporting evidence from the existence of the
cosmic microwave background and its explanation of the cosmic abundance of helium,
it should be kept in mind that none of its assumptions is unchallengeable. Many of
the calculations and much of the observational evidence which go into constructing
this history of the universe are reviewed by Peebles (1971) and Weinberg (1972), and
we refer the reader to these references for more quantitative details.

A good deal of the nature of the early universe can be understood from the fact that
the decrease of the scale factor \(a\) as one goes back toward the past has the same
local effect on the matter as if the matter were placed in a box whose walls contract
at the same rate. Thus, as would happen in a contracting box, the contribution of
radiation compared with ordinary matter (i.e., baryons) increases in the past, as was
already noted in Section~5.2. The energy density of the cosmic microwave background
in the present universe is estimated to be about 1,000 times smaller than the mass
density contribution of matter, although this number is uncertain by as much as a
factor of 10 because of the difficulties involved in determining the present density of
ordinary matter. Thus, assuming that this radiation continues to exist as one goes back
into the past (e.g., that it was not, say, emitted by galaxies), then according to
Eqs.~\eqref{eq:5.2.19} and \eqref{eq:5.2.20}, when the scale factor \(a\) was more
than 1,000 times smaller than its present value, this radiation should have been the
dominant contribution to the energy density of the universe. Thus, one would expect
the radiation filled model of the universe to be a good approximation for the dynamics
of the universe before this stage, while the dust filled model should be a good
approximation afterwards.

Just as the temperature of a box of gas increases as one compresses it, one would
expect the matter and radiation in the universe to get hotter as \(a\) decreases, and to
become infinitely hot as one approaches the big bang, \(a\to0\).\footnote{It would
also be self-consistent to assume that no radiation was present in the very early
universe and that only cold matter, such as baryons, was present. However, a ``cold
big bang'' model would have the major tasks of accounting for the present existence of
the cosmic microwave background and the correct helium abundance of the universe,
both of which are naturally explained by the standard ``hot big bang'' model described
here.} If the early universe was radiation dominated, as we expect it to be, then for
all models (\(k=0,\pm1\)) the dependence of \(a\) and \(\rho\) on \(\tau\) for small
\(\tau\) goes over to the \(k=0\) solution,
\begin{align}
  a(\tau)&=(4C')^{1/4}\tau^{1/2}, &&\tag{5.4.1}\label{eq:5.4.1}\\
  \rho&=\frac{3}{32\pi G\tau^2}, &&\tag{5.4.2}\label{eq:5.4.2}
\end{align}
where here, and throughout this section, we shall restore the \(G\)'s and \(c\)'s in
our equations. If the radiation is thermally distributed, the mass density \(\rho\) is
given by the following expression, derived from the quantum statistical mechanics of
massless particles:
\begin{equation}
  \rho=\sum_{i=1}^{n}\alpha_i g_i\frac{\pi^2}{30\hbar^3c^5}(kT)^4,
  \tag{5.4.3}\label{eq:5.4.3}
\end{equation}
where \(n\) is the number of species of radiation, \(g_i\) is the spin degeneracy
factor, and \(\alpha_i\) takes the value 1 for bosons and \(7/8\) for fermions.
Massive particles whose mass is much less than \(kT\) effectively act like zero mass
particles and should be included in Eq.~\eqref{eq:5.4.3} as a ``species of radiation.''
Equations~\eqref{eq:5.4.1}, \eqref{eq:5.4.2}, and \eqref{eq:5.4.3} imply that
\(T\propto\rho^{1/4}\propto a^{-1}\), as expected from the redshift relation,
Eq.~\eqref{eq:5.3.6}.

An important issue is whether the interactions of matter and radiation in the early
universe proceed on a rapid enough time scale for thermalization to occur locally
(i.e., within the particle horizon). If they do not, then the self-consistency of
assuming that the matter is thermally distributed is questionable, and the predicted
evolution of the matter and radiation may depend on the details of the assumed initial
distribution. If they do, then we have the relatively simple task of evolving a thermal
distribution of matter until such time as equilibrium is no longer maintained. The
expansion time scale, \(t_E\), of the universe, i.e., the time over which a significant
change in the scale factor \(a\) occurs, is
\begin{equation}
  t_E\sim\frac{a}{\dot a}=2\tau
  \tag{5.4.4}\label{eq:5.4.4}
\end{equation}
by Eq.~\eqref{eq:5.4.1}. On the other hand, the time scale for interactions is
\begin{equation}
  t_I\sim\frac{1}{n\sigma c}\propto\frac{a^3}{\sigma}
  \propto\frac{\tau^{3/2}}{\sigma(T)},
  \tag{5.4.5}\label{eq:5.4.5}
\end{equation}
where the number of interacting particles is assumed to be conserved, so that their
number density, \(n\), scales as \(a^{-3}\), and the possible energy dependence of
the interaction cross section \(\sigma\) is explicitly indicated by writing \(\sigma\)
as a function of temperature. Comparison of Eqs.~\eqref{eq:5.4.4} and
\eqref{eq:5.4.5} shows that unless \(\sigma\) falls off rapidly at high energies, at
sufficiently early times we will have \(t_I\ll t_E\), so thermalization can be achieved.
(Actually, it is possible that at very high energies the interactions of particle physics
become ``asymptotically free,'' and \(\sigma\) does fall off sufficiently rapidly to make
\(t_I>t_E\) as \(\tau\to0\). However, even if this occurs, for energies smaller than
about \(10^{15}\,\mathrm{GeV}\) we should recover \(t_I>t_E\), and thermalization still
should be achieved at a very early time.) On the other hand, as the universe evolves,
we eventually find \(t_I>t_E\), and the matter distribution will drop out of thermal
equilibrium.

Thus, the above considerations lead to the following basic picture of the evolution
of our universe. The universe began as a hot (\(T\to\infty\)), dense
(\(\rho\to\infty\)) ``soup'' of matter and radiation in thermal equilibrium. However,
as the universe evolved, thermal equilibrium was not maintained, as \(t_I\) became
greater than \(t_E\). The energy content of the early universe was dominated by
radiation. However, by the time \(a\) reached about \(1/1000\) of its present value,
the ordinary matter contribution dominated the energy content of the universe, and
the dynamics of the universe became that of a dust filled Robertson--Walker model.
We shall now proceed to fill in some of the important details of this evolutionary
history.

During the first \(10^{-43}\) seconds of the evolutionary history of the universe
predicted by classical general relativity, the magnitude of the curvature of spacetime
was greater than the scale given by the Planck length
\((G\hbar/c^3)^{1/2}\approx10^{-33}\,\mathrm{cm}\). On dimensional grounds, we
expect that the quantum effects of the gravitational field will be very important in
this era (see Chapter~14), and thus the predictions of classical general relativity
cannot be taken seriously. After this time, we may expect classical general relativity
to be valid, but at the extreme conditions found in the very early universe (at
\(\tau=10^{-43}\,\mathrm{s}\), we have \(\rho\sim10^{92}\,\mathrm{g\,cm^{-3}}\)!) it
hardly needs to be emphasized that all statements about the behavior of matter during
this epoch are speculative.

There are two interesting and important effects that may have occurred in the very
early universe, only several orders of magnitude later than the Planck time. The first
concerns the dynamics of the very early universe. Some quantum field theory models
which attempt to unify the strong and electro-weak interactions predict that at very
high temperatures, the thermal equilibrium state of the quantum field will undergo a
phase transition. In these models, if ``supercooling'' occurs there may be an important
``vacuum'' contribution, of the form \(-\lambda g_{ab}\) (where \(\lambda\) is a large,
positive constant) to the stress-energy tensor, \(T_{ab}\), of the field. Thus, the very
early universe may have gone through a phase where the dynamics was the same as
would occur in an empty universe with a large, positive cosmological constant. Hence,
if these models are correct, then, as first suggested by Guth (1981), there may be an
``inflationary'' regime in the very early universe where the universe is approximately
a de Sitter solution [see problem~3(b)] and expands very rapidly. If this occurs, then
as mentioned at the end of Section~5.3, the particle horizons in the present universe
could be much larger than would be calculated by extrapolating the solutions of
Table~\ref{tab:5.1} all the way back to the ``big bang.'' For further discussion of
inflationary cosmological models, we refer the reader to Gibbons, Hawking, and
Siklos (1983).

The second effect concerns the production of baryons. There is strong reason to
believe (Steigman 1976) that the matter content of the universe consists of baryons
as opposed to antibaryons; that one does not have matter-antimatter symmetry. It is
possible that our universe was simply born with an excess of baryons over antibaryons,
and, indeed this must be so if baryon number is strictly conserved. However, it is
also possible that the universe began in a matter-antimatter symmetric state and that
the baryon excess was manufactured in the very early universe. In order for this to
happen, it is necessary that the high energy particle interactions occurring in the very
early universe satisfy the following properties: (1) Clearly, they must fail to conserve
baryon number. (2) They must fail to preserve charge conjugation, \(C\), and the
composition of charge conjugation with parity, \(CP\). (If either of these symmetries
are preserved, equal numbers of baryons and antibaryons will be produced.) (3) They
must result in departures from thermal equilibrium. (This is because particles and
antiparticles have equal masses, so in thermal equilibrium they will occur in equal
numbers.) Nonequilibrium phenomena could be produced naturally by the existence of
a massive particle whose decay lifetime is greater than the expansion time, \(t_E\), at
the time when the temperature of the universe drops below the mass of the particle
(so that production of the particle is essentially shut off). Remarkably, the ``grand
unified'' theories of the strong, electromagnetic, and weak interactions currently
predict all three of these properties and thus may provide an explanation of the
matter-antimatter asymmetry in our universe.

We pick up our account of the history of the early universe at time \(\tau=1\) second,
when the density is \(\rho\approx5\times10^5\,\mathrm{g\,cm^{-3}}\), and the
temperature is \(T\approx10^{10}\,\mathrm{K}\). Although these conditions are extreme
by ordinary standards, we are now in a low enough energy and density regime to be
able to make solid predictions. At this time the matter in the universe consists almost
entirely of neutrinos, photons, electrons, positrons, neutrons, and protons in thermal
equilibrium; the temperature is low enough that the equilibrium abundance of the more
massive elementary particles is negligible. By about this stage, the interactions of the
neutrinos have become sufficiently weak that they decouple from the rest of the matter.
For the remainder of the history of universe, they merely passively get redshifted to
lower energy. Since a redshifted thermal spectrum (\(\omega\to\omega/a\)) is simply a
thermal spectrum at a lower temperature (\(T\to T/a\)), the present universe should
be filled with a blackbody distribution of neutrinos at temperature \(T\approx2\,\mathrm{K}\).
However, the detection of these neutrinos would be a task very far beyond presently
available sensitivities.

As the universe continues to cool, the rates of reactions which convert protons to
neutrons and vice versa quickly drop to much lower than the expansion rate of the
universe. Consequently, the neutron to proton ratio ``freezes out'' at
\(\tau\approx1.5\) seconds at roughly the value \(1/6\). (The protons are more
abundant than the neutrons because they are over \(1\,\mathrm{MeV}\) lighter and thus
are more prevalent in thermal equilibrium before the ``freezeout'' occurs.) Of course,
the term ``freezeout'' should not be taken too literally, since the ``turnoff'' of the
interactions does not occur instantaneously and, furthermore, the neutron-proton ratio
continues to decrease slowly with time due to the decay of the neutrons.

At \(\tau=4\) seconds, we have \(\rho=3\times10^4\,\mathrm{g\,cm^{-3}}\), and
\(T\approx5\times10^9\,\mathrm{K}\approx0.5\,\mathrm{MeV}\), which is approximately
the mass of electrons and positrons. At this stage the equilibrium population of
electrons and positrons decreases rapidly; the production rate drops below the
annihilation rate, and shortly after this time all the positrons will have annihilated,
leaving a relatively small population of residual electrons. Essentially all the energy
of the electron-positron pairs is transferred to the photons, heating them to a
temperature \(\sim1.4\) times higher than the temperature of the neutrinos.

When the temperature drops to about \(10^9\,\mathrm{K}\) at \(\tau\approx3\) minutes,
nucleosynthesis begins rather abruptly, producing \({}^4\mathrm{He}\) nuclei. Actually,
in thermal equilibrium the abundance of \({}^4\mathrm{He}\) at these baryon densities
would be appreciable at even higher temperatures (\(\sim3.5\times10^9\,\mathrm{K}\)),
but very little nucleosynthesis occurs before \(\tau=3\) minutes because deuterium
(\({}^2\mathrm{H}\)) plays a key role in the nuclear reactions which build up helium,
but the equilibrium abundance of \({}^2\mathrm{H}\) is very low until the temperature
drops to \(10^9\,\mathrm{K}\). Almost no nucleosynthesis of elements beyond
\({}^4\mathrm{He}\) occurs because of the large Coulomb barriers and the lack of stable
nuclei with atomic weights 5 and 8. Within a time span of a few minutes, essentially
all of the neutrons present at the ``freezeout'' time which did not decay are converted
to \({}^4\mathrm{He}\), resulting in an abundance of \({}^4\mathrm{He}\) of about 25\% by
mass, with much smaller abundances of \({}^2\mathrm{H}\), \({}^3\mathrm{He}\), and
\({}^7\mathrm{Li}\) also produced, but negligible amounts of other elements. The
percentage of \({}^4\mathrm{He}\) is not very sensitive to the assumed value of baryon
density, since it is governed mainly by the neutron-proton ratio at ``freezeout,'' but
the abundances of the other elements, particularly \({}^2\mathrm{H}\), are very sensitive
to the baryon density. A relatively low baryon density can produce a \({}^2\mathrm{H}\)
abundance of over \(5\times10^{-4}\) by mass, whereas a high baryon density---
specifically, a density high enough to ``close the universe,'' as discussed below---
increases the efficiency of the chain of reactions producing \({}^4\mathrm{He}\), and
results in a \({}^2\mathrm{H}\) abundance many orders of magnitude lower.

It is difficult to observe the cosmic abundance of helium, but the abundance of 25\%
seems to be in agreement with observations. The presence of this amount of helium in
the universe cannot readily be accounted for by other processes---in particular,
nucleosynthesis in stars is estimated to produce an abundance of helium of only a few
percent---and thus the prediction of helium production via ``big bang nucleosynthesis''
must be viewed as a major success of the theory.

After the period of nucleosynthesis, the universe, of course, continued its expansion
and cooling. The next cosmic occurrence of major importance took place when the
temperature had dropped to about \(4000\,\mathrm{K}\). This occurred at
\(\tau\sim4\times10^5\) years if the universe was still radiation dominated, but it
would have occurred somewhat earlier if the universe had already become matter
dominated. At this temperature and below it, the free electrons and protons combined
to form neutral hydrogen. Indeed, by the time the temperature had dropped to
\(2000\,\mathrm{K}\), the fraction of ionized hydrogen was only \(\sim10^{-4}\). As a
result of this occurrence---called \emph{recombination}, although, of course, the
electrons and protons had never combined earlier---the interaction between the matter
and radiation dropped precipitously, since the scattering cross section of photons off
free charged particles is much greater than off neutral \(\mathrm{H}\) and \(\mathrm{He}\).
Indeed, the photons effectively decoupled completely from the matter after
recombination, and merely cooled with the expansion of the universe during the
remainder of its evolution. Thus, the present universe should be filled with this low
temperature blackbody radiation originating from the big bang, whose photons last
interacted with matter at the recombination time.

Exactly such a radiation background at temperature \(T\approx2.7\,\mathrm{K}\)
(corresponding to wavelengths mainly in the microwave regime) was discovered by
Penzias and Wilson (1965). The existence of this radiation would be difficult to
account for in any other way, and thus it provides a major confirmation of the above
picture of the evolution of our universe. Furthermore, the radiation has been measured
to be isotropic to a high degree of precision. (A ``dipole type'' anisotropy of
\(\sim0.1\%\) has been detected [Smoot \emph{et al.}\ 1977], but this presumably is
due to motion of the Earth with respect to the ``preferred rest frame'' of the
Robertson--Walker cosmology.) This provides very strong evidence that the universe
was very nearly homogeneous and isotropic at least down to the recombination time.

The decoupling of matter and radiation in this era also had a major effect on the growth
of gravitational perturbations, leading to the formation of galaxies. Just before
recombination, the pressure provided by the radiation inhibited the growth of
gravitational perturbations involving masses smaller than about \(10^{17}M_\odot\),
which is much larger than typical galactic masses (\(\sim10^{11}M_\odot\)). However,
after recombination, the radiation pressure had no effect on the matter, and
gravitational instabilities could occur on all mass scales \(\gtrsim10^5M_\odot\). Thus,
irregularities in the distribution of matter began to grow after recombination,
resulting in the formation of galaxies, star clusters, and stars. However, the details of
this process are not very well understood at present (see Peebles 1980 for further
discussion).

At some time between \(\tau\sim10^3\) years and \(\tau\sim10^7\) years, ordinary
matter became the dominant form of energy in the universe. (The exact time at which
this occurred is uncertain because there is substantial uncertainty in the matter density
of the present universe.) The dynamics of the universe was transformed from a
``radiation filled'' to a ``dust filled'' solution. Finally, at \(\tau\sim10\text{--}20\)
billion years the universe reached its present state.

Thus, general relativity, together with the assumption of homogeneity and isotropy
and assumptions about the matter content of the universe, has produced a remarkably
successful picture of the history of our universe. Among its most notable achievements
are its explanations of the cosmic abundance of helium and the existence of the cosmic
microwave radiation.

It is interesting that if one accepts this picture, some nontrivial constraints are placed
on (1) the masses of stable, weakly interacting, elementary particles and (2) the number
of species of massless particles which are in thermal equilibrium in the early universe.
To explain the first constraint, suppose that the electron or muon neutrino has a mass,
\(m\), or that some massive stable particle exists which participates only in the weak
interactions. Then the behavior of this particle in the early universe would be the same
as that of the massless neutrino. However, in the present universe, instead of
contributing energy \(\sim10^{-4}\,\mathrm{eV}\) per particle (corresponding to
\(T\sim2\,\mathrm{K}\)), it would contribute energy \(m\) per particle. If the particles
are very massive, their population will be ``frozen in'' at a very low value when they
drop out of thermal equilibrium in the early universe and they will not contribute much
to the present energy density of the universe. But if their mass lies within the range
\(10\,\mathrm{GeV}\gtrsim m\gtrsim100\,\mathrm{eV}\), these particles would become by
far the dominant contributors of mass-energy in the universe, and they would produce
dynamical effects on the expansion rate of the universe which are incompatible with
observations. Thus, our cosmological theory and observations place some stringent
limits on the possible masses of stable, weakly interacting particles.

The existence of other species of massless particles (e.g., neutrinos) in thermal
equilibrium in the early universe would not significantly affect the present mass density
of the universe, since their energy would get redshifted away as the universe expands.
However, they would have an important dynamical effect upon the early universe. This
is because, according to Eq.~\eqref{eq:5.4.3}, they would affect the relation between
\(\rho\) and \(T\), making \(T\) smaller for a given \(\rho\). Since for small \(\tau\)
the relation between \(\rho\) and \(\tau\) always is given by Eq.~\eqref{eq:5.4.2} for a
radiation dominated universe, we see that \(T(\tau)\) will be decreased, i.e., a given
temperature will occur at an earlier epoch, when the expansion rate is higher. The most
important consequence of this effective speedup of the expansion rate of the universe
is that the neutron to proton ratio will ``freeze out'' at a higher temperature, yielding a
higher percentage of neutrons. Consequently, more helium will be manufactured during
the era of nucleosynthesis.\footnote{However, if the expansion rate is tremendously
speeded up, there will not even be enough time for the nucleosynthetic reactions and
the helium synthesis will decrease.} Assuming that the baryon density in the early
universe corresponds at least to the present mass density observed in small groups of
galaxies (see below), it is found that an unacceptably high abundance of helium will
be manufactured if there are more than four species of neutrinos. Since the electron
and muon neutrinos are known to exist and the existence of a neutrino associated with
the recently discovered \(\tau\)-lepton seems likely, this cosmological limit appears
to leave little room for the existence of other species of massless particles in thermal
equilibrium in the early universe.

What about the future evolution of our universe? The most important issue with
regard to future evolution is whether our universe is ``open'' or ``closed,'' i.e., does
it correspond to the cases \(k=0,-1\), or the case \(k=+1\)? As discussed in
Section~5.2, if the universe is open it will expand forever, while if it is closed it
will eventually recontract. We may express the basic equations
Eqs.~\eqref{eq:5.2.14} and \eqref{eq:5.2.15} governing the dynamics of the present
universe in terms of Hubble's constant, \(H=\dot a/a\), and the \emph{deceleration
parameter}, \(q\), defined by
\begin{equation}
  q=-\frac{\ddot a\,a}{(\dot a)^2}.
  \tag{5.4.6}\label{eq:5.4.6}
\end{equation}
Since \(P\simeq0\) in the present universe, we have
\begin{align}
  H^2&=\frac{8\pi G\rho}{3}-\frac{kc^2}{a^2}, &&\tag{5.4.7}\label{eq:5.4.7}\\
  q&=\frac{4\pi G\rho}{3H^2}. &&\tag{5.4.8}\label{eq:5.4.8}
\end{align}
Defining \(\Omega\) by
\begin{equation}
  \Omega=\frac{8\pi G\rho}{3H^2},
  \tag{5.4.9}\label{eq:5.4.9}
\end{equation}
we see that \(q=\Omega/2\), and the universe is closed (\(k=1\)) if and only if
\(\Omega>1\), i.e., \(\rho>\rho_c\equiv3H^2/8\pi G\).

There are, at present, four basic, independent pieces of observational evidence, which,
in conjunction with the above equations, yield indications of whether the universe is
open or closed: (i) the redshift--apparent magnitude relation for distant objects, (ii)
the mass density of the present universe, (iii) the age of the universe, and (iv) the
cosmic abundance of deuterium. We shall briefly discuss each of these in turn.

As derived at the end of Section~5.3a, for sufficiently nearby objects we have a
linear relation between redshift, \(z\), and present proper distance, \(R\),
Eq.~\eqref{eq:5.3.9}. Also, for sufficiently nearby objects, the apparent luminosity of
an object is proportional to its intrinsic luminosity divided by \(R^2\). (For very
distant objects, one must correct this simple relationship because of the expansion of
the universe and the curvature of space.) Thus, if we consider objects of a fixed
intrinsic luminosity, for sufficiently nearby objects, one will obtain a linear relation
between \(\ln z\) and \(\ln[\text{apparent luminosity}]\) = apparent magnitude. We
will not discuss here the difficult tasks of (a) how to find ``standard candles,'' i.e.,
objects of fixed intrinsic luminosity, and (b) how to calibrate their intrinsic luminosity
(or, equivalently, their distance) by a chain of arguments whereby nearby standard
candles are used to calibrate more distant standard candles (see e.g., Weinberg 1972).
The relation between \(\ln z\) and apparent magnitude is, indeed, found to be linear,
corresponding to a value of \(H\) given by
\begin{equation}
  H\sim50\,\mathrm{km\,s^{-1}\,Mpc^{-1}},
  \tag{5.4.10}\label{eq:5.4.10}
\end{equation}
where \(1\) megaparsec (Mpc) \(\approx3\times10^{19}\,\mathrm{km}\). Because of
uncertainties in the absolute distance determinations, this value of \(H\) quite
possibly could be in error by as much as a factor of 2 and, indeed, some recent
determinations have given the value \(H\sim100\,\mathrm{km\,s^{-1}\,Mpc^{-1}}\). For
sufficiently distant objects the relation between apparent magnitude and \(\ln z\) will
depart from linearity in a way that depends on the dynamical evolution of \(a(\tau)\)
and, to a lesser extent, on the spatial geometry. In principle, by examining these
departures from linearity, we can determine \(q\), which, in turn, will tell us if the
universe is open (\(q\leq1/2\)) or closed (\(q>1/2\)). However, in practice, these
departures begin to occur only for objects at such large distances
(\(z\gtrsim0.2\), corresponding to \(R\gtrsim10^3\,\mathrm{Mpc}\)) that the light
reaching us was emitted at a significantly earlier epoch. Thus, there is reason to doubt
whether our most distant standard candles---the brightest galaxies in clusters of
galaxies---are really ``standard'' on account of evolutionary effects; these galaxies may
well have been significantly brighter (or dimmer) in the past. Until these evolutionary
effects are better understood, it will not be possible to draw definitive conclusions as
to whether the universe is open or closed by this approach.

The parameter \(\Omega\) can be determined directly by measuring the mass density,
\(\rho\), of the present universe and Hubble's constant \(H\). As mentioned above,
the uncertainty in calibration of distances leads to considerable uncertainty in \(H\).
Fortunately, however, the distance scale calibration also enters into the determination
of \(\rho\) in such a way that the ratio \(\Omega\propto\rho/H^2\) is unaffected by
these uncertainties, and thus \(\Omega\) may be determined much more precisely than
either \(\rho\) or \(H\). In the present universe, ordinary matter appears to make the
dominant contribution to \(\rho\); the energy density of the cosmic microwave
background is \(\sim10^{-3}\) times the estimated matter density quoted at the end of
this paragraph. It appears that essentially all the matter in the universe is concentrated
in galaxies, since most possibilities for intergalactic matter are excluded on
observational or theoretical grounds. The masses of spiral galaxies sufficiently nearby
(\(\lesssim15\,\mathrm{Mpc}\)) can be estimated from their rotational velocities. The
mass of a cluster of galaxies can be estimated from the random velocities of the
individual galaxies, assuming that the cluster is gravitationally bound. When these
masses are compared, some interesting discrepancies arise: The masses of galaxies in
binary systems and small clusters as determined by the second method appears to be
higher (by a factor of about 3) than the masses of galaxies determined by the first
method. Furthermore, the masses of galaxies in large clusters appears to be larger (by
a factor of at least 2 and possibly much larger) than the masses obtained for small
clusters. A possible explanation of the first discrepancy is that galaxies possess a
massive halo of underluminous material; this extra mass would not affect the galactic
rotation curve and would cause the first method (which estimates only the mass of the
luminous inner portion of the galaxy) to underestimate the total mass. The cause of
the second discrepancy is less clear, but it might be explained by the presence of
intercluster material or simply by the fact that galaxies in large clusters might tend to
be significantly more massive than galaxies in small clusters. If we take the typical
mass of galaxies to be that determined by the second method for small clusters, we
obtain \(\Omega\sim0.04\), corresponding to \(\rho\sim2\times10^{-31}\,
\mathrm{g\,cm^{-3}}\) if the value of \(H\) given in Eq.~\eqref{eq:5.4.10} is used.
This provides evidence in favor of an open universe, but many possibilities for hidden
mass remain, so this evidence is not conclusive.

Since our universe is presently matter dominated and should have been so for almost
all of its history, from our solution of Table~\ref{tab:5.1} we find that if \(k=0\) the
age of the universe is related to Hubble's constant by
\begin{equation}
  \tau_c=\frac{2}{3H}\sim1.3\times10^{10}\ \text{years},
  \tag{5.4.11}\label{eq:5.4.11}
\end{equation}
using the value of \(H\) given in Eq.~\eqref{eq:5.4.10}. If \(k=-1\), \(\tau\) will be
larger than \(\tau_c\), while if \(k=+1\), it will be smaller. Thus, comparison of the
actual age of the universe with the critical age, \(\tau_c\), will tell us whether the
universe is open or closed. Recall, however, that there is considerable uncertainty in
the value of \(H\), and thus in the value of \(\tau_c\). The age of the universe can be
estimated from the age of the oldest objects we can find: globular star clusters (whose
age can be estimated from the theory of stellar evolution) and elements manufactured
by the first stars (whose age can be determined by radioactive dating if the element is
radioactive). Study of these objects yields an age of the universe in the range of 10 to
20 billion years. The rough equality of this age with \(\tau_c\) provides further strong
confirmation of the ``big bang'' origin of the universe. However, the uncertainties in
both are too great to allow us to conclude that the universe is open or that it is closed.

A final clue as to whether the universe is open or closed is provided by the cosmic
abundance of deuterium. As mentioned above, if the density of baryons in the early
universe was sufficient to give a baryon density in the present universe greater than
the critical density \(\rho_c=3H^2/8\pi G\), then very little deuterium should remain
after the era of nucleosynthesis. However, measurements of the cosmic abundance of
deuterium have yielded mass fractions of \(\sim2\times10^{-5}\), corresponding to a
present value of \(\Omega\) of \(\sim0.1\). This provides evidence that the universe is
open, but significant loopholes remain, such as the possibility that deuterium was
manufactured by some other process, or that most of the present mass density of the
universe is in nonbaryonic matter (e.g., a neutrino species with mass \(\sim100\,
\mathrm{eV}\)).

In summary, many important issues in cosmology remain to be resolved. But it is
already clear that general relativity has provided us with a remarkably successful
picture of the spacetime structure of our universe.
